\chapter{The model in brief}\label{ch:model}

This chapter describes the stochastic actor-oriented model (SAOM) and its estimation in the terms that \nwsaom{} uses. It is not a substitute for the literature~\cite{snijders2010,ripley}; its purpose is to define the quantities that appear in the output: the rates, the coefficients, the standard errors and the convergence $t$-ratios.

\section{Data and the idea of the model}

The data are a network observed at waves $t_1<t_2<\dots<t_W$ on a fixed set of $n$ actors. Let $x(t_m)$ be the adjacency matrix at wave $m$, with $x_{ij}=1$ if actor $i$ has a tie to actor $j$. The observations are viewed as snapshots of a continuous-time Markov process $X(t)$ on the set of all networks: between two waves ties change one at a time, in the order of unobserved \emph{ministeps}. Each ministep is the decision of one actor to change one of its own outgoing ties (or to leave things as they are). The process is simulated between waves $m$ and $m+1$ starting from the observed $x(t_m)$ over a time interval of length one; the \emph{rate} (below) therefore has the meaning of the expected number of opportunities for change per actor between two waves, and each period $m=1,\dots,W-1$ has its own rate $\rho_m$.

\section{The rate function}

In each period every actor $i$ gets opportunities to make a ministep at a rate $\lambda_i$. In the basic model all actors share the rate of the period, $\lambda_i=\rho_m$. The waiting time until the next opportunity of anyone is exponentially distributed with parameter $\sum_i\lambda_i$, and the actor who gets the opportunity is $i$ with probability $\lambda_i/\sum_h\lambda_h$. With \opt{ratecov()} the rate depends on actor covariates $z_{ih}$,
\begin{equation}
  \lambda_i = \rho_m\,\exp\Bigl(\sum_h b_h z_{ih}\Bigr),
\end{equation}
and the coefficients $b_h$ are estimated jointly with the other parameters (they are shared by all periods; each period keeps its own $\rho_m$). The structural rate effects (\opt{outrate} and relatives) let the rate depend on the actor's current degree, so that it changes during the simulation.

\section{The objective function and the ministep}

When actor $i$ gets an opportunity, it chooses among the $n$ alternatives of changing the tie variable $x_{ij}$ for one $j\neq i$ (creating the tie if it does not exist, withdrawing it if it does) and of changing nothing. Write $x(i\leadsto j)$ for the network that results from the alternative $j$, with $x(i\leadsto i)=x$. The actor evaluates each alternative with an \emph{objective function} (in RSiena's terminology the \emph{evaluation function}),
\begin{equation}
  f_i(\beta,x)=\sum_k \beta_k\, s_{ik}(x),
\end{equation}
a weighted sum of effect statistics $s_{ik}$, and chooses with the multinomial-logit probability
\begin{equation}
  P(i \text{ changes } x_{ij}\mid x) = \frac{\exp\bigl(f_i(\beta,x(i\leadsto j))\bigr)}{\sum_{h}\exp\bigl(f_i(\beta,x(i\leadsto h))\bigr)} .
  \label{eq:choice}
\end{equation}
The statistics $s_{ik}$ describe the local network around $i$. In \nwsaom{} they are selected by options; for example the \emph{outdegree} (density) effect has $s_{i}(x)=\sum_j x_{ij}$, the \emph{reciprocity} effect $s_i(x)=\sum_j x_{ij}x_{ji}$ and the \emph{transitive triplets} effect (\opt{transtrip}) counts, in RSiena's definition, the triples $(i,j,h)$ with $x_{ij}=x_{ih}=x_{hj}=1$. A positive coefficient means that an actor is, other things equal, more likely to choose alternatives that increase the statistic. The coefficients are log-odds in the sense of~\eqref{eq:choice}; the estimates for outdegree are typically negative because ties are sparse.

The model is \emph{actor-oriented} and \emph{myopic}: an actor evaluates the alternatives only through its own statistics $s_{ik}$, not through the effect that its choice would have on other actors' statistics. This is the property that distinguishes a SAOM from an exponential random graph model, which has no actors and no ministeps, but one probability distribution over whole networks.

\paragraph{Endowment and creation.} The objective function can be split by the direction of the change. A \emph{creation} effect contributes only when a tie is created, an \emph{endowment} effect only when a tie is withdrawn (a tie, once it exists, is ``endowed'' with extra value). The options \opt{endow()} and \opt{creation()} select such effects; the statistics used in the estimation are, for the endowment role, minus the sum over the ties lost between two waves of the effect's tie statistic in the start network and, for the creation role, the sum over the ties gained of the tie statistic in the end network.

\section{Behavior and several networks}

A behavior variable $z_i$ with a bounded range of integer values can evolve together with the network. It has its own rate function and its own objective function. A behavior ministep has three alternatives, to decrease the value by one, to keep it, or to increase it by one (clamped at the range of the observed values), and they are chosen by the same multinomial-logit rule~\eqref{eq:choice}. Opportunities for network ministeps and behavior ministeps come from competing processes: one waiting time is drawn from the total rate summed over all variables, then the variable that acts is drawn in proportion to its share of the total rate, and then the actor. Selection (network effects that depend on the behavior) and influence (behavior effects that depend on the network) are thus estimated in one model. In a model with several dependent networks, \texttt{nwsaom multinet}, the network that changes is drawn in the same way.

\section{Estimation by the method of moments}

Let $\theta=(\beta,\rho)$ be the vector of all parameters and $S(X)$ the vector of the statistics that correspond to the parameters. The method of moments chooses $\theta$ such that the expected value of the statistics at the end of the simulated period, given the observed start, equals what was observed,
\begin{equation}
  E_\theta\bigl[S(X(t_{m+1}))\mid X(t_m)=x(t_m)\bigr]=S(x(t_{m+1})) ,
  \label{eq:mom}
\end{equation}
summed over the $W-1$ periods (the effect statistics are pooled over periods, each period keeping its own rate). The expectation cannot be computed in closed form, so it is replaced by simulation, and the equation is solved by stochastic approximation (Robbins--Monro). \nwsaom{} follows the three phases of RSiena's \texttt{siena07}.
\begin{description}
  \item[Phase 1] $k_0$ simulations at the starting values (option \opt{k0()}, default 50) estimate the derivative matrix $D=\partial E_\theta S/\partial\theta$. It is estimated with the score-function method as the covariance of the statistics with the scores of the simulated ministeps; the phase ends with a first step towards the solution.
  \item[Phase 2] Iterations $\theta^{(N+1)}=\theta^{(N)}-a_N\,D^{-1}\bigl(S^{(N)}-s_{\mathrm{obs}}\bigr)$, where $S^{(N)}$ is the statistic of one simulated period and $s_{\mathrm{obs}}$ the observed one. The gain $a_N$ is reduced over four subphases, starting with \opt{firstg()} (default 0.2) and multiplied by 0.5 in each next subphase, as RSiena's defaults.
  \item[Phase 3] $k_3$ simulations at the final estimate (option \opt{k3()}, default 1\,000) estimate $D$ and the covariance matrix $\Sigma$ of the statistics again. The covariance matrix of the estimates is
  \begin{equation}
    V = D^{-1}\,\Sigma\,(D^{-1})' ,
  \end{equation}
  which is \texttt{e(V)}; the standard errors in the table are its diagonal elements' square roots. The same simulations give the convergence statistics.
\end{description}
The first and the third phase are the part of the algorithm that is trivially parallel, and \nwsaom{} runs their replicates on worker threads (Section~\ref{sec:threads}).

\section{Conditional and unconditional estimation}

The rate parameters can be handled in two ways.

\paragraph{Unconditional estimation.} Each rate $\rho_m$ is a parameter of the method of moments like the others. The statistic that corresponds to a rate is the \emph{distance}, the number of dyads in which the network at the end of the unit-length period differs from the network at its start; its observed value is the number of dyads that differ between the two waves. The rates are estimated together with the effects, with standard errors that come from the same sandwich formula. Option \opt{unconditional} requests this for a network-only model.

\paragraph{Conditional estimation.} The rates are not Robbins--Monro parameters. Each period is simulated until the distance from the start equals the observed distance between the two waves, so that the number of changes is conditioned on, and the effect statistics are those of the network at that point. A period's rate estimate is the mean of the time the simulation needed, over the $k_3$ phase-3 replicates, and its reported standard error is the standard deviation of these times. There are no rate $t$-ratios. Conditional estimation is the default for a model with one dependent network and no behavior, as in RSiena.

Co-evolution models, models with several networks and network-only models in which composition change restricts which actors are present are always estimated unconditionally. A rate covariate coefficient stays a method-of-moments parameter under conditional estimation. The two estimators are consistent for the same model but differ in finite samples; in RSiena's own tests on the Glasgow data the estimates differed by up to $0.3$ standard errors.

\section{Convergence}\label{sec:convergence}

After phase 3, for each parameter the \emph{convergence $t$-ratio} is computed as the mean deviation $S_k^{(r)}-s_{\mathrm{obs},k}$ over the $k_3$ simulated replicates divided by its standard deviation. If the estimate solves~\eqref{eq:mom}, the mean deviation is zero. As in RSiena, the fit is considered converged when every $t$-ratio is below $0.1$ in absolute value and the \emph{overall maximum convergence ratio} (RSiena's \texttt{max convergence ratio}, the largest convergence $t$-ratio over linear combinations of the parameters; stored in \texttt{e(tconv\_max)}) is below $0.25$. The ratios are printed below the coefficient table for all parameters, rates included, and stored in \texttt{e(tconv)}; \texttt{e(tratio)} holds those of the coefficients in \texttt{e(b)}.

These $t$-ratios are not the $z$-statistics of the table, which are the coefficients divided by their standard errors, and a large convergence ratio does not mean a significant effect: it means that the algorithm did not reach the solution. The usual remedy is to run the model again, starting from the estimates (options \opt{theta0()} and, for unconditional fits, \opt{rate0()}); an example is in Chapter~\ref{ch:start}. A model that cannot be made to converge, or that stops with the message that the estimation diverged, is usually too rich for the data. Parameters that are not identified because there are too few instances of an effect in the data (isolate effects on small networks, endowment and creation effects of rare events) are a frequent cause.

\section{Non-directed and two-mode networks}

In a non-directed network every tie is symmetric ($x_{ij}=x_{ji}$), so that a tie is the joint property of two actors and the model has to say who decides. \nwsaom{} offers RSiena's model types, selected with \opt{symtype()}: in the \emph{forcing} model (the default) an actor chooses a tie to create or end and imposes the change on the other actor; in the \emph{confirmation} model a new tie also needs the other actor's consent; and in the pairwise models \emph{force}, \emph{agree} and \emph{joint} a pair of actors meets and the tie changes if the initiator wants it, if both want a new tie or either wants to end one, or if their summed preferences favour the change. In a two-mode network the actors are the nodes of one mode, and they choose ties to the nodes of the other mode (the events), which do not act. See the chapter on network types.
